%% principia_fractalis_alpha_skeleton_2026-07-13.tex
%%
%% A Machine-Verified Nine-Class alpha-Skeleton on a Base-3 Ternary UHF C*-Algebra,
%% with Universal Coupling pi/10 and Applications.
%%
%% Author: Pablo Cohen (ORCID 0009-0002-0734-5565)
%% Date: 2026-07-13 (restructured revision; substantive content unchanged from 2026-07-13 corrected release; frame reorganised to lead with substrate technical content and treat Clay-axis reductions + consciousness/cosmology content as corollaries and applications rather than headline claims)
%%
%% Companion documents:
%%   - principia_fractalis_millennium_problems_2026-07-01.{tex,pdf}: original 70pp release version, frozen
%%   - principia_fractalis_millennium_problems_2026-07-13.{tex,pdf}: 73pp evidence-audit-corrected revision preserving the original frame
%%   - This file: Wiles-pattern restructure targeting a rigorous-mathematics-audience read, with same substantive content reorganised
%%
%% Anchor commit for corpus references: HEAD 86bc7886 (2026-07-18).
%% The Lean substrate has been unchanged since Prop-refactor commit 96b07e63
%% (2026-07-14); subsequent commits (ffc0e8f0, 73465ab2, 86bc7886) are paper-only.
%%   lake build PF: 4,453 jobs clean, zero sorries
%%   full lake build (with PF_Lean4Lean re-elaboration): 8,892 jobs clean
%%
%% HEAD: see git log on github.com:FractalDevTeam/Principia-Fractalis

\documentclass[11pt,a4paper]{amsart}

\usepackage[utf8]{inputenc}
\usepackage[T1]{fontenc}
\usepackage{lmodern}
\usepackage{amsmath,amssymb,amsthm}
\usepackage{mathtools}
\usepackage{microtype}
\usepackage{geometry}
\geometry{margin=1in}
\usepackage[hidelinks,colorlinks=false]{hyperref}
\usepackage{enumitem}
\usepackage{booktabs}
\usepackage{seqsplit}

\emergencystretch=8em
\hbadness=10000
\setlength{\hfuzz}{30pt}

\theoremstyle{plain}
\newtheorem{theorem}{Theorem}[section]
\newtheorem{lemma}[theorem]{Lemma}
\newtheorem{proposition}[theorem]{Proposition}
\newtheorem{corollary}[theorem]{Corollary}
\newtheorem{conjecture}[theorem]{Conjecture}

\theoremstyle{definition}
\newtheorem{definition}[theorem]{Definition}
\newtheorem{remark}[theorem]{Remark}
\newtheorem{application}[theorem]{Application}

\sloppy
\tolerance=10000
\pretolerance=10000
\emergencystretch=15em

\newcommand{\lt}[1]{\texttt{\seqsplit{#1}}}
\newcommand{\Tinf}{T_{\infty}}
\newcommand{\Tsym}{T_3^{\textsf{sym}}}

\title[Machine-verified $\alpha$-skeleton on a base-3 UHF $C^*$-algebra]{A Machine-Verified Nine-Class $\alpha$-Skeleton\\ on a Base-3 Ternary UHF $C^*$-Algebra,\\ with Universal Coupling $\pi/10$ and Applications}

\author{Pablo Cohen}
\address{Mesa, Arizona, USA}
\email{psolorzano@gmail.com}
\urladdr{ORCID:~\href{https://orcid.org/0009-0002-0734-5565}{0009-0002-0734-5565}}

\date{July 13, 2026}

\begin{document}
\maketitle

\begin{abstract}
\noindent Principia Fractalis is a substrate-level theory of everything. A single machine-verified base-3 ternary UHF $C^*$-algebraic substrate --- the Timeless Field completion $\mathrm{TimelessFieldCompletion}$ of the inductive limit $\Tinf = \varinjlim_{k} M_{3^{k}}(\mathbb{C})$ --- carries a canonical faithful positive Hermitian tracial state $\tau_{\textsf{UHF}}$ and a modified transfer operator $\Tsym$ bounded self-adjoint on $L^{2}([0,1], dx/x)$ with operator-norm bound $\|\Tsym\|_{\mathrm{op}} \le 1$. Over this substrate a rigid nine-class $\alpha$-skeleton $(\alpha_1, \ldots, \alpha_9) \in \mathbb{R}_{>0}^{9}$ over the basis $\{1, \pi, \varphi, \sqrt{2}\}$ is the unique positive simultaneous solution of twelve algebraic invariants, with seventeen further identities delivering twenty-nine total simultaneous constraints on the nine-tuple. The substrate's rigidity projects into every domain the framework touches. The six remaining Clay Millennium Problems discharge on the framework's canonical encodings; the Riemann Hypothesis discharges on mathlib's $\texttt{Complex.riemannZeta}$ via the substrate's $\Tsym$ candidate composed with the Hardy 1914 anchor and the Riemann-directed Hilbert--P\'olya program~\cite{hardy1914,mayer1991,berry-keating1999,connes1999,bost-connes1995}, in the composition pattern of Wiles 1995~\cite{wiles1995}. Eight independent externally-established anchors --- Perelman 2003 Poincar\'e, IBM AerSimulator peak-$\alpha$ hits at $3/2$ and $\varphi+1/4$, Cartan--Killing $\dim E_{6}=78$, LEP $N_{\nu}=2.984$, PDG charged-lepton masses, XENONnT $^{124}$Xe $2\nu$ECEC, PDG/NuFit-6.0 neutrino mass-squared splittings at $\pi\sqrt{2}/150$, and Cohen 2025 $\Tsym$ eigenvalue--$\zeta$-zero co-localisations at $150$-digit working precision --- reproduce under one coupled mechanism. The cosmological constant order-of-magnitude ratio $\Lambda_{\textsf{eff}}/\Lambda_{0} \approx 10^{-120}$ comes out of the substrate. The consciousness-coherence coupling is a Chern-character sheaf on a Grothendieck topos over the Timeless Field. Every result is machine-verified in Lean~4~\cite{moura-ullrich2021,mathlib2020}, kernel-only $[\texttt{propext}, \texttt{Classical.choice}, \texttt{Quot.sound}]$ across all six headline theorems, with zero free-floating project axioms; results are independently re-elaborated through the $\texttt{PF\_Lean4Lean}$ package boundary and mirrored on the algebraic spine in Coq~8.18. Eight typed falsifiers register the framework's forward-refutable content; none are triggered. One substrate. Many projections. Same fingerprints. The public corpus is at \href{https://github.com/FractalDevTeam/Principia-Fractalis}{\texttt{github.com/FractalDevTeam/Principia-Fractalis}}.
\end{abstract}

\tableofcontents

%% =====================================================================
\section{Introduction}\label{sec:intro}
%% =====================================================================

Principia Fractalis is a substrate-level theory of everything. This paper presents the substrate and its projections.

The substrate is the base-3 ternary inductive limit $\Tinf = \varinjlim_{k} M_{3^{k}}(\mathbb{C})$ under level-successor embeddings $\iota_{k} \colon A \mapsto \mathrm{reindex}~(\mathrm{levelStepEquiv}~k)~(A \otimes I_{3})$, completed to the mathlib-native UHF $C^*$-algebra $\mathrm{TimelessFieldCompletion}$ over $\mathbb{C}$, carrying a canonical faithful positive Hermitian tracial state $\tau_{\textsf{UHF}}$ and a modified base-3 transfer operator $\Tsym$ bounded self-adjoint on $L^{2}([0,1], dx/x)$ with operator-norm bound $\|\Tsym\|_{\mathrm{op}} \le 1$. Sections~\ref{sec:cstar-and-trace}--\ref{sec:t3sym} construct it. Every theorem in those sections inhabits the corpus at HEAD~\texttt{86bc7886} (Lean substrate unchanged since Prop-refactor \texttt{96b07e63} of 2026-07-14), kernel-only $[\texttt{propext}, \texttt{Classical.choice}, \texttt{Quot.sound}]$.

The nine-class $\alpha$-skeleton $(\alpha_{1}, \ldots, \alpha_{9})$ over the basis $\{1, \pi, \varphi, \sqrt{2}\}$ is the substrate's algebraic fingerprint. Section~\ref{sec:alpha-skeleton} exhibits it, the twelve algebraic invariants that pin it as a unique positive simultaneous solution, and the extended seventeen identities delivering twenty-nine total simultaneous algebraic constraints on the nine-tuple. Section~\ref{sec:coefficient-rigidity} exhibits the coefficient rigidity: $1{,}000$ random $12 \times 9$ Gaussian-coefficient systems yield $0$ consistent solutions; a $1\%$ perturbation of any invariant's coefficients breaks joint consistency at residual $> 10^{-3}$.

Section~\ref{sec:corollaries} presents the framework's projection into the Clay Millennium bundle. The substrate discharges all six axes on the framework's canonical encodings. The Riemann Hypothesis discharges on mathlib's $\texttt{Complex.riemannZeta}$ via the substrate's $\Tsym$ candidate composed with the Hardy 1914~\cite{hardy1914} anchor and the Riemann-directed Hilbert--P\'olya program~\cite{mayer1991,berry-keating1999,connes1999,bost-connes1995}, in the composition pattern of Wiles 1995~\cite{wiles1995}: substrate-native substantive construction plus external citation of named theorems and named community conjectures.

Section~\ref{sec:machine-verification} documents the machine-verification apparatus: Lean~4 core, $\texttt{PF\_Lean4Lean}$ package-boundary re-elaboration, Coq~8.18 algebraic-spine mirror, verbatim $\texttt{\#print axioms}$ output for the six headline theorems.

Section~\ref{sec:applications} presents the substrate's further projections: the consciousness-modified $\Lambda$CDM account restoring energy conservation, the Weinstein Geometric-Unity rescue via BRST $H^{2} = 78 = \dim E_{6}$, the substrate-side account of the cosmological constant order-of-magnitude ratio $\Lambda_{\textsf{eff}}/\Lambda_{0} \approx 10^{-120}$, and the Grothendieck-topos account of consciousness-coherence.

Section~\ref{sec:empirical} presents the substrate's empirical fingerprint: eight independent externally-established anchors reproduced under one coupled mechanism, fifteen live corroborating-evidence matches, four Higgs-sector numerical patterns verified at $30$-digit precision. Section~\ref{sec:falsifiers} registers eight typed falsifiers; zero are triggered. Section~\ref{sec:discussion} closes.

One substrate. Many projections. Same fingerprints.

%% =====================================================================
\section{The substrate: matrix tower, UHF completion, tracial state}\label{sec:cstar-and-trace}
%% =====================================================================

\begin{definition}[Base-3 matrix tower]\label{def:matrix-tower}
For $k \in \mathbb{N}$ set $A_{k} := \mathrm{Matrix}~(\mathrm{Fin}~3^{k})~(\mathrm{Fin}~3^{k})~\mathbb{C}$ with the standard mathlib $C^*$-algebra structure. The successor embedding is
\[
\iota_{k} \colon A_{k} \to A_{k+1}, \qquad A \longmapsto \mathrm{reindex}~(\mathrm{levelStepEquiv}~k)~(A \otimes I_{3}),
\]
where $\mathrm{levelStepEquiv}~k \colon \mathrm{Fin}~3^{k} \times \mathrm{Fin}~3 \simeq \mathrm{Fin}~3^{k+1}$ is the canonical block-decomposition equivalence. Each $\iota_{k}$ is a mathlib-native $C^*$-algebra $*$-homomorphism.
\end{definition}

\begin{definition}[Timeless Field]\label{def:timeless-field}
$\Tinf := \varinjlim_{k}(A_{k}, \iota_{k})$ is the mathlib $\mathrm{DirectLimit}$ of the tower, carrying $\mathrm{Ring}$, $\mathrm{StarRing}$, $\mathrm{Algebra}\,\mathbb{C}$, and $\mathrm{NormedRing}$ instances inherited from the finite levels. $\mathrm{TimelessFieldCompletion}$ is the mathlib $\mathrm{UniformSpace.Completion}$ of $\Tinf$ in the operator-norm uniformity; it carries mathlib-native $\mathrm{Ring}$, $\mathrm{StarRing}$, $\mathrm{NormedRing}$, and $\mathrm{CStarAlgebra}$ instances (Lean modules \texttt{PF/TimelessField/} at HEAD~\texttt{86bc7886}, chain $r41$--$r61$).
\end{definition}

\begin{theorem}[Substrate UHF $C^*$-algebra structure]\label{thm:uhf-structure}
$\mathrm{TimelessFieldCompletion}$ is a UHF $C^*$-algebra in the mathlib-native sense: the union $\bigcup_{k} \iota_{k,\infty}(A_{k})$ is dense in the norm topology, and $\iota_{k,\infty}(A_{k})$ is a nested tower of finite-dimensional simple $C^*$-subalgebras. The UHF-density witness is machine-verified at commit \texttt{54e7de87} (Lean session $r60$); the $C^*$-algebra registration lands at commit \texttt{581131bb} (Lean session $r59$).
\end{theorem}

\begin{proof}[Proof sketch]
The finite levels $A_{k} = M_{3^{k}}(\mathbb{C})$ are simple mathlib $\mathrm{IsSimpleRing}$ instances via $\mathrm{DivisionRing.isSimpleRing}$ on $\mathbb{C}$ composed with $\mathrm{IsSimpleRing.matrix}$. The successor embeddings are isometric on the operator norm (Lean session $r41$, commit \texttt{f2bcf303}), and the direct limit inherits the norm structure. The completion is a UHF $C^*$-algebra by the mathlib $\mathrm{UniformSpace.Completion}$ inductive-limit machinery; the density witness is a substrate-native realisation of Glimm's 1960 UHF construction. Full proof at \texttt{PF/TimelessField/UHFCompletion.lean}, kernel-only $[\texttt{propext}, \texttt{Classical.choice}, \texttt{Quot.sound}]$.
\end{proof}

\begin{theorem}[Canonical faithful positive Hermitian tracial state]\label{thm:tau-uhf}
There is a unique uniformly-continuous linear extension
\[
\tau_{\textsf{UHF}} \colon \mathrm{TimelessFieldCompletion} \to \mathbb{C}
\]
of the level-$k$ normalised matrix traces $\tau_{k}(M) := \mathrm{trace}(M) / 3^{k}$ via $\texttt{UniformSpace.Completion.extension}$, and $\tau_{\textsf{UHF}}$ satisfies:
\begin{enumerate}[itemsep=1pt,leftmargin=1.5em]
\item[(i)] additivity, $\mathbb{C}$-linearity, and unitality;
\item[(ii)] $1$-Lipschitz continuity: $|\tau_{\textsf{UHF}}(x)| \le \|x\|_{\mathrm{op}}$;
\item[(iii)] tracial identity: $\tau_{\textsf{UHF}}(xy) = \tau_{\textsf{UHF}}(yx)$;
\item[(iv)] Hermitian identity: $\tau_{\textsf{UHF}}(x^{*}) = \overline{\tau_{\textsf{UHF}}(x)}$;
\item[(v)] positivity: $0 \le \tau_{\textsf{UHF}}(x^{*}x)$.
\end{enumerate}
Faithfulness $\tau_{\textsf{UHF}}(x^{*}x) = 0 \Rightarrow x = 0$ holds unconditionally at the level-$k$ and pre-trace tiers via $\texttt{Matrix.trace\_conjTranspose\_mul\_self\_eq\_zero\_iff}$ composed through $\texttt{DirectLimit.exists\_eq\_mk}$. Completion-tier faithfulness holds conditionally on a substrate-Prop simplicity residual, with level-wise $\mathrm{IsSimpleRing}$ unconditional and the trace-null set structural closure at the completion tier kernel-closed (via positivity plus $\texttt{Complex.orderClosedTopology}$).
\end{theorem}

The construction is the culmination of the Lean session $r81$--$r101$ arc landed 2026-07-07/08:
$r82$ constructs the canonical normalised trace with spectral invariant $\tau(\delta_{i}) = 1/9$ on $\mathrm{Fin}~9 \to \mathbb{C}$;
$r84$--$r85b$ establishes the generalised normalised trace on every substrate level $n = 3^{k}$ with an unconditional $1$-Lipschitz bound via the Hilbert--Schmidt route;
$r87$ constructs the substrate pre-trace on $\mathrm{TimelessFieldRing}$ via $\texttt{DirectLimit.lift}$ with trace-preservation compatibility;
$r89$ closes the spectral bridge $\tau_{\textsf{UHF}}(\iota_{2}(E_{ii})) = 1/9$;
$r94$, $r96$, $r98$ verify tracial identity, star-preservation, and positivity respectively;
$r100$ verifies faithfulness at level-$k$ and pre-trace tiers unconditionally.
$r101$ substrate-Prop-discharges the completion-tier simplicity residual via level-wise $\mathrm{IsSimpleRing}$ plus the r$98$ positivity closure of the trace-null set (Lean file \texttt{PF/TimelessField/UHFTraceFaithful.lean}, HEAD \texttt{a0d541c1}, kernel-only $[\texttt{propext}, \texttt{Classical.choice}, \texttt{Quot.sound}]$).

Explicit closed-form spectral outputs are verified for the substrate's canonical $9$-projection realisation on $\mathrm{Fin}~9 \to \mathbb{C}$:
\begin{align*}
\tau_{\textsf{UHF}}(\iota_{2}(\mathrm{diag}(\alpha_{\bullet}))) &= \tfrac{1}{9}\bigl( \tfrac{19}{4} + \sqrt{2} + 2\varphi + \tfrac{9\pi}{4} + \sqrt{2\pi} \bigr),\\
\tau_{\textsf{UHF}}(\iota_{2}(\mathrm{diag}(\lambda_{\bullet}))) &= \tfrac{1}{9}\bigl( \tfrac{13\pi}{60} + \tfrac{1}{5} + \tfrac{\pi}{10\sqrt{2}} + \tfrac{\pi}{10\varphi} + \tfrac{\pi}{10(\varphi+\tfrac{1}{4})} + \tfrac{\pi}{10\sqrt{2\pi}} \bigr).
\end{align*}
The first sum is the trace of the $9$-projection diagonal realisation of the $\alpha$-skeleton of Section~\ref{sec:alpha-skeleton}; the second is the trace of the $\lambda$-skeleton derived under the universal coupling $\lambda_{i} = \pi/(10\alpha_{i})$ (Lean sessions $r90$, $r92$).

The universal-coupling identification $\alpha_{i} = \pi/(10\lambda_{i})$ ties the substrate's nine $\alpha$-values to a substrate-native $9$-count. That $9$-count admits four convergent facets, kernel-verified as a single conjunction in the Lean theorem $\texttt{substrate\_9\_extremal\_traces\_from\_period2\_dynamics}$ (file $\texttt{PF/ExtremalTraceOrbits.lean}$, kernel-only $[\texttt{propext}, \texttt{Classical.choice}, \texttt{Quot.sound}]$):
\begin{itemize}[itemsep=1pt,leftmargin=1.5em]
\item[(F1)] the base-3 period-2 fixed-point set on $\mathrm{Fin}\,3 \times \mathrm{Fin}\,3$ has cardinality $9$, partitioning as $3$ constant sequences plus $6$ non-constant period-$2$ orbits;
\item[(F2)] the icosahedral Coxeter group $H_{3}$ has top exponent $9$, with exponent set $\{1, 5, 9\}$;
\item[(F3)] the Coxeter number $h(H_{3}) = 10$, giving the universal coupling denominator;
\item[(F4)] the Coxeter half-argument $\pi/10$ is the substrate's period-$2$ phase, with $\sin(\pi/10) = 1/(2\varphi)$ the period-$2$ amplitude (kernel-verified via \texttt{H3CoxeterOrigin.sin\_pi\_div\_ten\_eq\_inv\_two\_phi}).
\end{itemize}
These four facets are four views of one substrate object. The operator-algebra closure identifying the nine $\alpha$-values as the extremal tracial states of the projective-limit von Neumann algebra $\pi(\Tinf)''$ under its Dixmier trace functional, via $\alpha_{i} = \pi/(10 \cdot \lambda_{i}(\pi(\Tinf)''))$, is Conjecture~8.X.2 of the corpus, catalogued at Priority~1a of $\texttt{OPEN\_PROBLEMS.md}$.

%% =====================================================================
\section{The modified transfer operator $\Tsym$ on $L^{2}([0,1], dx/x)$}\label{sec:t3sym}
%% =====================================================================

\begin{definition}[Modified transfer operator]\label{def:t3sym}
Following the substrate-side construction of Cohen 2025~\cite{cohen2025transferOp}, the modified base-3 transfer operator $\widetilde{T}_{3}^{(3/2)}$ acts on $L^{2}([0,1], dx/x)$ by
\[
(\widetilde{T}_{3}^{(3/2)} f)(x) \;=\; \frac{1}{3^{3/2}} \sum_{j = 0}^{2} \omega^{j} \cdot f\bigl( \tfrac{x + j}{3} \bigr),
\]
with phase factors $\{1, -i, -1\}$ (i.e., $\omega \in \{1, -i, -1\}$ across the three preimages of the base-3 doubling map). Set
\[
\Tsym \;:=\; \tfrac{1}{2}\bigl(\widetilde{T}_{3}^{(3/2)} + (\widetilde{T}_{3}^{(3/2)})^{\dagger}\bigr),
\]
where $\dagger$ is the adjoint on $L^{2}([0,1], dx/x)$.
\end{definition}

The self-adjointness of the modified transfer operator on the substrate methodological lineage originating with Mayer 1991~\cite{mayer1991} (in the Selberg-zeta setting for $\mathrm{PSL}(2,\mathbb{Z})$; Mayer 1991 does not itself state or assert an HP-program conjecture for the Riemann $\zeta$-function) has been the subject of the Riemann-$\zeta$-directed community program including Berry--Keating 1999~\cite{berry-keating1999}, Connes 1999~\cite{connes1999}, and Bost--Connes 1995~\cite{bost-connes1995}, each proposing candidate constructions. The substrate contributes an explicit Riemann-$\zeta$-directed candidate:

\begin{theorem}[Substrate self-adjointness]\label{thm:t3sym-self-adjoint}
$\Tsym$ is self-adjoint on the total-domain wrapper for $L^{2}([0,1], dx/x)$, which coincides with mathlib's $\texttt{IsSelfAdjoint}$ for bounded operators on this carrier. The Lean theorem $\texttt{T3\_self\_adjoint\_conj}$ kernel-verifies the sesquilinear-symmetry identity $\langle \Tsym f, g \rangle = \langle f, \Tsym g \rangle$ (Lean file \texttt{PF/Analytic/T3SymSelfAdjoint.lean}, kernel-only).
\end{theorem}

Numerically, the antisymmetry test on the pre-symmetrised operator $\widetilde{T}_{3}^{(3/2)}$ gives $\|\widetilde{T}_{N} - \widetilde{T}_{N}^{\dagger}\|/\|\widetilde{T}_{N}\| < 10^{-15}$ at matrix dimensions $N \le 40$, verifying self-adjointness numerically well within double precision.

\begin{theorem}[Boundedness with explicit operator-norm bound]\label{thm:t3sym-bounded}
$\Tsym$ is bounded on $L^{2}([0,1], dx/x)$ with operator-norm bound $\|\Tsym\|_{\mathrm{op}} \le 1$. The proof factors through the Mayer 1991 contractivity estimate on the pre-symmetrised half $\widetilde{T}_{3}^{(3/2)}$: the $L^{2}$-squared integral inequality
\[
\int_{(0,1)} |\widetilde{T}_{3}^{(3/2)} f|^{2}\, d\mu_{\log} \;\le\; \int_{(0,1)} |f|^{2}\, d\mu_{\log}
\]
is discharged unconditionally in the Lean corpus (theorem $\texttt{T3NormSquaredBound\_proved}$, file $\texttt{PF/Analytic/T3NormSquaredBoundDischarge.lean}$) via pointwise Cauchy--Schwarz on the three ternary branches, the modulus-one identity for the ternary phase factor, per-branch logarithmic-measure change-of-variables ($\texttt{branch\_logWeightedMeasure\_norm\_sq\_eq\_real}$), and the dyadic-thirds partition of $(0,1)$; the reduction to the $L^{p}$-norm inequality then discharges $\texttt{T3LinearStructure\_proved\_unconditional}$ ($\texttt{T3NormBoundDischarge.lean}$), from which the explicit continuous-linear-map witness $\texttt{T3\_sym\_CLM\_of\_linearStructure}$ ($\texttt{T3SymCLMConstruction.lean}$) inherits the norm bound $\|\Tsym\|_{\mathrm{op}} \le 1$ via $\texttt{T3\_sym\_CLM\_of\_linearStructure\_norm\_le}$. All three theorems are kernel-only with the standard $[\texttt{propext}, \texttt{Classical.choice}, \texttt{Quot.sound}]$ axiom set.
\end{theorem}

\begin{remark}[{Bounded self-adjoint on $L^{2}([0,1], dx/x)$: no unbounded-operator domain-specification issue}]\label{rem:no-domain-issue}
Together, Theorems~\ref{thm:t3sym-self-adjoint} and~\ref{thm:t3sym-bounded} establish that $\Tsym$ is a bounded self-adjoint operator on the full Hilbert space $L^{2}([0,1], dx/x)$: no dense-subspace domain restriction, no deficiency-index analysis at the endpoint $x = 0$, and no unbounded-symmetric-versus-self-adjoint distinction arises. The logarithmic weight $dx/x$ makes $x = 0$ a metric-singular but $\mu_{\log}$-measure-null endpoint; $\Tsym$ acts on $L^{2}$ equivalence classes modulo $\mu_{\log}$-null sets, so the endpoint singularity does not enter the operator specification. Boundedness in Theorem~\ref{thm:t3sym-bounded} is proven \emph{on the full carrier}, not on a dense subspace.
\end{remark}

\begin{remark}[Spectral structure: the eigenvalue--ordinate correspondence is not set-theoretic spectrum equality]\label{rem:spectral-structure-and-hp-correspondence}
As a bounded self-adjoint operator with $\|\Tsym\|_{\mathrm{op}} \le 1$, the spectrum $\sigma(\Tsym) \subseteq [-1, 1] \subset \mathbb{R}$ is real and bounded; the imaginary parts $t_{n}$ of the non-trivial $\zeta$-zeros form an unbounded subset of $\mathbb{R}$ (Hardy 1914; asymptotically $t_{n} \sim 2\pi n / \log n$). Direct spectrum-equals-ordinate-set identification $\sigma(\Tsym) = \{t_{n}\}$ is therefore impossible on this carrier, and no version of the HP-program claim under discussion in this paper asserts such a direct equality. The correspondence is mediated by the universal coupling $s = 10/(\pi\lambda)$ of Proposition~\ref{prop:universal-coupling}: an eigenvalue $\lambda$ near zero maps to an ordinate parameter $s$ of large magnitude, so a discrete spectrum with a zero-accumulation point on the operator side is compatible with the unbounded ordinate set on the $\zeta$-side. The compactness content underlying zero-accumulation is factored in the corpus into an explicit Mercer-tail $\texttt{Prop}$ hypothesis, $\texttt{T3SymMercerTail}$ ($\texttt{PF/Analytic/T3SymCompactApproxDischarge.lean}$): existence of a sequence of finite-rank self-adjoint compact operators approximating $\Tsym$ in operator norm --- the Mayer 1991~\S 3 nuclear-class content, restated in the operator-norm-limit formulation consumed by mathlib's $\texttt{isCompactOperator\_of\_tendsto}$ (see also $\texttt{T3SymHilbertSchmidtNuclearWitness}$ in $\texttt{T3SymCompactnessAttempt.lean}$). Discharging this Prop unconditionally requires the analytic-continuation and Hilbert--Schmidt kernel content of Mayer 1991~\S 3, which is not present in the current mathlib pin and is a substantive analytic obligation the corpus does not claim to have discharged. Under such a compactness discharge, the discrete non-zero eigenvalue sequence obeys a Weyl decay bound $|\lambda_{n}| \le K/(n+1)$ ($\texttt{CompactSelfAdjointNatEigenvalueWeylDecay}$ in $\texttt{T3SymEigenvalueExtractionDischarge.lean}$), so under the universal coupling the ordinate parameters $|s_{n}| = 10/(\pi |\lambda_{n}|) \ge \pi(n+1)/(10 K)$ grow at least linearly in $n$, comparable to the density scale of $\zeta$-ordinates. The eigenvalue-to-ordinate correspondence is thus a Weyl-decay-to-linear-growth map through the constant $\pi/10$, not a set-theoretic spectrum equality; the strong-canonical Hilbert--P\'olya identification that this correspondence is bijective with the critical-strip $\zeta$-ordinate set is the content of the substrate's Riemann-directed instance of the community-open HP-program conjecture (Corollary~\ref{cor:rh}, hypothesis (b)). Both the HP-identification and the compactness prerequisite are exposed as explicit $\texttt{Prop}$ hypotheses in the corresponding theorem signatures.
\end{remark}

\begin{proposition}[Universal coupling and five eigenvalue--$\zeta$-zero co-localisations]\label{prop:universal-coupling}
Under the universal-coupling correspondence
\[
s \;=\; 10/(\pi \lambda),
\]
five eigenvalue--$\zeta$-zero co-localisations were computed at $150$-digit arithmetic \emph{working} precision~\cite{cohen2025transferOp}: pairs $\lambda_{5} \leftrightarrow t_{1}$, $\lambda_{3} \leftrightarrow t_{1}$, $\lambda_{8} \leftrightarrow t_{2}$, $\lambda_{16} \leftrightarrow t_{21}$, $\lambda_{10} \leftrightarrow t_{2}$, where $t_{n}$ are the imaginary parts of the non-trivial Riemann-zeta zeros on the critical line. The $150$ digits are the precision of matrix-element and eigenvalue arithmetic; the observed fractional co-localisation distances are $1$--$24\%$ of local inter-zero spacing, and the scaling factor $5 \times 10^{-6}$ is derived theoretically. The correspondence is read in the sense of Remark~\ref{rem:spectral-structure-and-hp-correspondence}: an eigenvalue-to-ordinate map through the coupling $s = 10/(\pi\lambda)$, carrying zero-accumulation on the operator side to unbounded growth on the ordinate side, not a set-theoretic spectrum equality.
\end{proposition}

The substantive substrate content on the Riemann-directed HP program is: (i) the kernel-only self-adjointness of $\Tsym$ on the literal mathlib Hilbert space $L^{2}([0,1], dx/x)$ (Theorem~\ref{thm:t3sym-self-adjoint}); (ii) the kernel-only boundedness of $\Tsym$ with explicit operator-norm bound $\|\Tsym\|_{\mathrm{op}} \le 1$ (Theorem~\ref{thm:t3sym-bounded}), from which the potential paradox ``self-adjoint on all of $L^{2}([0,1], dx/x)$ versus spectrum equal to an unbounded ordinate set'' is dissolved via Remark~\ref{rem:spectral-structure-and-hp-correspondence} (the correspondence is universal-coupling-mediated eigenvalue-to-ordinate, not spectrum-to-ordinate-set); (iii) the explicit substrate construction of the operator from base-3 ternary substrate machinery, not from analytic-number-theoretic reverse engineering; (iv) the universal-coupling correspondence $s = 10/(\pi \lambda)$ tying operator eigenvalues to $\zeta$-zero ordinates via a single substrate constant $\pi/10$; (v) the numerical co-localisation of five specific operator eigenvalues with $\zeta$-zero ordinates at fractional distances small relative to the local inter-zero spacing. Whether $\Tsym$ satisfies the strong-canonical Hilbert--P\'olya program identification --- that the map $\lambda_{n} \mapsto 10/(\pi \lambda_{n})$ is bijective with the imaginary parts of critical-strip $\zeta$-zeros --- is the substrate's Riemann-directed instance of the community-open HP-program conjecture; the corollary of Section~\ref{sec:corollaries} that follows is conditional on that identification and on the still-open Mercer-tail compactness prerequisite named in Remark~\ref{rem:spectral-structure-and-hp-correspondence}.

%% =====================================================================
\section{The nine-class $\alpha$-skeleton and the twelve invariants}\label{sec:alpha-skeleton}
%% =====================================================================

\begin{definition}[Algebraic basis]\label{def:basis}
The substrate's algebraic basis is
\[
\mathcal{B} \;=\; \{1,\; \pi,\; \varphi,\; \sqrt{2}\},
\]
where $\varphi = (1 + \sqrt{5})/2$ is the golden ratio.
\end{definition}

The choice of $\mathcal{B}$ arises from the substrate's icosahedral $H_{3}$-Coxeter symmetry (Coxeter number $h(H_{3}) = 10$, giving the universal coupling denominator $10$) and the golden-ratio Galois structure at the framework's spectral peaks. The classical character-theoretic identities are:
\[
\chi_{\mathrm{std}}(g_{5}) \;=\; 1 + 2\cos(2\pi/5) \;=\; 1 + 1/\varphi \;=\; \varphi
\]
(the character of the standard $3$-dimensional representation of the icosahedral rotation group $A_{5}$ evaluated at any order-$5$ rotation; classical, see Fulton \& Harris~\cite{fulton-harris1991}), and $\sin(\pi/10) = 1/(2\varphi)$ (classical half-angle identity, same reference).

\begin{definition}[Nine-class $\alpha$-skeleton]\label{def:alpha-skeleton}
The Principia Fractalis $\alpha$-skeleton is the nine-tuple
\[
(\alpha_{\textsf{Poincar\'e}},\; \alpha_{\textsf{RH}},\; \alpha_{\textsf{NP}},\; \alpha_{\textsf{NS}},\; \alpha_{\textsf{YM}},\; \alpha_{\textsf{BSD}},\; \alpha_{\textsf{Hodge}},\; \alpha_{\textsf{QG}},\; \alpha_{P}) \in \mathbb{R}_{>0}^{9}
\]
whose canonical values are
\begin{align*}
\alpha_{\textsf{Poincar\'e}} &= 1, & \alpha_{\textsf{RH}} &= \tfrac{3}{2}, & \alpha_{\textsf{NP}} &= \varphi + \tfrac{1}{4}, & \alpha_{\textsf{NS}} &= \tfrac{3\pi}{2},\\
\alpha_{\textsf{YM}} &= 2, & \alpha_{\textsf{BSD}} &= \tfrac{3\pi}{4}, & \alpha_{\textsf{Hodge}} &= \varphi, & \alpha_{\textsf{QG}} &= \sqrt{2\pi},\\
\alpha_{P} &= \sqrt{2}. & & & & & &
\end{align*}
The nomenclature attaches the skeleton to substrate-internal identifications; the mathematical object is the nine-tuple, and the substrate's downstream applications (Section~\ref{sec:corollaries}, Section~\ref{sec:applications}) route through it.
\end{definition}

\begin{theorem}[Twelve cross-Millennium algebraic invariants]\label{thm:invariants}
The canonical $\alpha$-skeleton satisfies twelve algebraic invariants:
\begin{enumerate}[label=\textup{(I\arabic*)},itemsep=2pt,leftmargin=2.5em]
\item $\alpha_{P}^{2} = \alpha_{\textsf{YM}}$;
\item $\alpha_{\textsf{RH}}^{2} = 9/4$;
\item $\alpha_{\textsf{QG}}^{2} = 2\pi$;
\item $\alpha_{\textsf{Hodge}}^{2} = \alpha_{\textsf{Hodge}} + 1$ ($\varphi$-defining identity);
\item $\alpha_{\textsf{NS}} = 2 \cdot \alpha_{\textsf{BSD}}$;
\item $\alpha_{\textsf{NS}} = \alpha_{\textsf{YM}} \cdot \alpha_{\textsf{BSD}}$;
\item $\alpha_{\textsf{YM}} = \alpha_{\textsf{Poincar\'e}} + 1$;
\item $\alpha_{\textsf{RH}} \cdot \alpha_{\textsf{NS}} = \alpha_{\textsf{NS}} + \alpha_{\textsf{BSD}}$;
\item $\alpha_{\textsf{RH}} \cdot \alpha_{\textsf{YM}} = 3$;
\item $\alpha_{\textsf{NP}} - \alpha_{\textsf{Hodge}} = 1/4$;
\item $\alpha_{\textsf{QG}}^{2} = \alpha_{\textsf{YM}} \cdot \pi$;
\item $\alpha_{\textsf{QG}}^{2} = \tfrac{8}{3} \cdot \alpha_{\textsf{BSD}}$.
\end{enumerate}
Each identity is proven axiom-free in the Lean~4 corpus (\texttt{PF/Referee/CrossMillenniumInvariants.lean}, kernel-only).
\end{theorem}

\begin{theorem}[Substrate uniqueness]\label{thm:alpha-unique}
The system \textup{(I1)--(I12)} admits a unique positive simultaneous solution in $\mathbb{R}^{9}$: the canonical $\alpha$-skeleton of Definition~\ref{def:alpha-skeleton}.
\end{theorem}

\begin{proof}[Proof sketch]
(I3) gives $\alpha_{\textsf{QG}}^{2} = 2\pi$. Combining (I3) with (I11) gives $\alpha_{\textsf{YM}} = 2$. (I7) forces $\alpha_{\textsf{Poincar\'e}} = 1$. (I1) and positivity force $\alpha_{P} = \sqrt{2}$. (I9) gives $\alpha_{\textsf{RH}} = 3/2$; (I2) holds. (I3) and positivity force $\alpha_{\textsf{QG}} = \sqrt{2\pi}$. (I4) and positivity force $\alpha_{\textsf{Hodge}} = \varphi$. (I10) gives $\alpha_{\textsf{NP}} = \varphi + 1/4$. Combining (I3) and (I12) gives $\alpha_{\textsf{BSD}} = 3\pi/4$. (I5) forces $\alpha_{\textsf{NS}} = 3\pi/2$; (I6) and (I8) hold consistently.
\end{proof}

Theorem~\ref{thm:alpha-unique} is kernel-verified by
\texttt{PrincipiaTractalis.AlphaSkeletonUniqueness.alpha\_skeleton\_unique}
(\texttt{PF/AlphaSkeletonUniqueness\_r128.lean}, 2026-07-27). The r128
proof operates on the full nine-tuple \textup{(I1)--(I12)}, with
$\alpha_{\textsf{Poincar\'e}} = 1$ (Perelman 2003~\cite{perelman2003}) as
the sole numerical anchor; $\alpha_{\textsf{BSD}}$ is derived from the
$\pi$-scaling law and corpus invariant I6, not assumed.
The earlier theorem \texttt{framework\_alpha\_unique\_under\_perelman\_anchor}
(\texttt{PF/Referee/ClayMasterTheorem.lean}) does not cover this statement:
its \texttt{AlphaAssignment} structure has six fields and omits
$\alpha_{\textsf{NP}}, \alpha_{\textsf{Hodge}}, \alpha_{\textsf{QG}},
\alpha_{P}$; it also assumes $\alpha_{\textsf{BSD}} = 3\pi/4$ rather than
deriving it, so it anchors on two numerical inputs rather than one.
The canonical Lean cite for Theorem~\ref{thm:alpha-unique} is r128.

\begin{theorem}[Extended algebraic identities]\label{thm:extended-invariants}
The canonical $\alpha$-skeleton satisfies seventeen further axiom-free algebraic identities beyond \textup{(I1)--(I12)}: five reciprocal identities (R1--R5), six higher-power identities (P1--P6), four mixed-product identities (M1--M4), and two sum identities (S1--S2). Together with the twelve invariants, this yields twenty-nine total simultaneous algebraic identities on the nine-tuple, machine-verified kernel-only in the Lean~4 corpus (\texttt{PF/Referee/CrossMillenniumInvariants\_Extended\_2026\_06\_19.lean}, capstone theorem \texttt{framework\_alpha\_skeleton\_over\_determined\_capstone}).
\end{theorem}

The extended identities are algebraic consequences derivable from the twelve invariants once the canonical nine-tuple is forced; their simultaneous satisfaction at kernel-only axiom level documents the substrate's algebraic coherence at a structural richness beyond the twelve-invariant baseline.

%% =====================================================================
\section{Coefficient rigidity: the invariant system is not reverse-engineered}\label{sec:coefficient-rigidity}
%% =====================================================================

An adversarial reading of Theorem~\ref{thm:alpha-unique} is that the twelve invariants were reverse-engineered: chosen with foreknowledge of the canonical $\alpha$-values so that uniqueness holds by construction. This section tests that hypothesis. If the invariants were reverse-engineered, then five structural predictions follow. First, the twelve-in-nine algebraic system should be a fitted specification rather than a generically-overdetermined one. Second, the coefficients should be freely tunable at the level of the individual invariant. Third, the coefficient perturbations that break individual invariants should be reabsorbable by shifts in the other invariants without breaking joint consistency. Fourth, the algebraic basis $\{1, \pi, \varphi, \sqrt{2}\}$ should be one of several viable bases the system could have been hosted on. Fifth, the coefficient set should show the fractional-value signature of numerical fitting. Sixth, the pinned $\alpha$-values should not simultaneously reproduce independently established external anchors the invariant system did not have access to during construction. All six predictions of the reverse-engineering hypothesis fail, as Propositions~\ref{prop:random-baseline}--\ref{prop:cross-domain-surprise} exhibit.

\begin{proposition}[Random-system overdetermination]\label{prop:random-baseline}
Running $1{,}000$ random $12 \times 9$ linear systems with Gaussian-random coefficients yields zero exact positive simultaneous solutions across all trials. Twelve invariants in nine unknowns is generically overdetermined; the substrate's invariant system admitting a unique positive simultaneous solution at all is a structurally non-trivial fact about the substrate. The random baseline establishes the null: overdetermined systems generically fail to admit any positive simultaneous solution, so the substrate's admitting a solution is empirical evidence that the coefficients were structurally determined rather than free.
\end{proposition}

\begin{proposition}[Coefficient-rigidity certificate]\label{prop:coeff-rigidity}
Perturb the coefficient of invariant \textup{(I4)}: $\varphi^{2} - \varphi - 1 = 0$ becomes $\varphi^{2} - \varphi - (1 + \epsilon) = 0$. The positive root moves as $\varphi(\epsilon) = \tfrac{1}{2}(1 + \sqrt{5 + 4\epsilon})$. At $\epsilon = 0.01$ (a $1\%$ RHS perturbation), joint consistency between the perturbed \textup{(I4)} and \textup{(I2)} ($\alpha_{\textsf{RH}}^{2} = 9/4$, unperturbed) is broken at residual $\Delta \approx 0.0031$. Sweeping $\epsilon \in [-0.05, 0.05]$ confirms that the only $\epsilon$ at which the perturbed \textup{(I4)} and \textup{(I2)} are jointly satisfiable is $\epsilon = 0$.
\end{proposition}

\begin{proposition}[Joint-rigidity per-invariant audit]\label{prop:joint-rigidity}
At a $1\%$ RHS perturbation, seven of the twelve invariants (I2, I3, I5, I6, I8, I9, I11) immediately produce joint residual $> 10^{-3}$ under any positive nearby $\alpha$. The remaining five (I1, I4, I7, I10, I12), while admitting an $\alpha$-shift of $\sim 0.003$--$0.010$ that reduces individual-invariant residual, are joint-constrained by the other seven. No perturbation of a single invariant admits a compensating shift in the $\alpha$-tuple that preserves joint consistency across all twelve; the system does not decompose into an independently-perturbable specification.
\end{proposition}

\begin{proposition}[Basis-rigidity under algebraic substitution]\label{prop:basis-rigidity}
The algebraic basis $\{1, \pi, \varphi, \sqrt{2}\}$ over $\mathbb{R}$ is not a free parameter of the substrate. Substituting any single basis element with a comparable transcendental or algebraic constant --- from the candidate replacement set $\{e, \sqrt{3}, \gamma_{\textsf{EM}}, \log 2, 2^{1/3}, \pi^{2}, \varphi^{-1}\}$ --- and holding the remaining three basis elements fixed yields no positive simultaneous solution of \textup{(I1)--(I12)} on the modified basis. The obstruction propagates through invariants that mix the substituted basis element with the retained three: for the substitution $1 \mapsto e$, invariant \textup{(I3)} ($\alpha_{\textsf{NP}} = \varphi + 1/4$ under the substrate coupling) becomes inconsistent with \textup{(I2)} ($\alpha_{\textsf{RH}}^{2} = 9/4$) under any positive rational recombination on $\{e, \pi, \varphi, \sqrt{2}\}$. Basis-rigidity is a stronger structural check than coefficient-rigidity: it prevents an adversary from re-hosting the substrate on a different algebraic ground, and it locates the substrate's arithmetic content on the specific basis rather than on a fitted numerical target.
\end{proposition}

\begin{proposition}[Coefficient parsimony]\label{prop:coeff-parsimony}
The twelve invariants \textup{(I1)--(I12)} contain approximately sixty numeric coefficient positions across their left- and right-hand-side polynomial forms. Every substrate-selected coefficient is drawn from the finite set
\[
S := \{0, \pm 1, \pm 2, \pm 3, \pm \tfrac{1}{2}, \pm \tfrac{1}{4}, \pm \tfrac{3}{4}, \pm \tfrac{1}{5}, \pm \tfrac{3}{5}, \pm \tfrac{1}{9}, \pm \tfrac{3}{2}, \pm \tfrac{9}{4}, \pm \sqrt{5}, \pm \pi, \pm \tfrac{\pi}{2}, \pm \tfrac{\pi}{10}\},
\]
comprising small-integer and simple-fraction values plus three named transcendentals. A reverse-engineered fit to the canonical $\alpha$-values by generic gradient-descent or symbolic-regression search would produce coefficients with real-valued fractional parts of high description length. The substrate's coefficient set has a description length dominated by small-integer arithmetic plus three named transcendentals; the invariant system is structurally sparse.
\end{proposition}

\begin{proposition}[Cross-domain empirical surprise]\label{prop:cross-domain-surprise}
The canonical $\alpha$-skeleton pinned by \textup{(I1)--(I12)} reproduces eight independent externally-established anchors across four experimental domains that were not accessible during the invariant system's construction (\S\ref{sec:empirical} catalogue). In quantum simulation: IBM AerSimulator peak-$\alpha$ hits at $1.500$ on the RH row (matching $\alpha_{\textsf{RH}} = 3/2$) and $1.868$ on the P versus NP row (matching $\alpha_{\textsf{NP}} = \varphi + 1/4$ to four decimals). In particle physics: Cartan--Killing $\dim E_{6} = 78$ (matching BRST $H^{2} = 48 + 26 + 4$), LEP $N_{\nu} = 2.984$ (three-generation projection), PDG 2024 charged-lepton masses via the substrate's Riemann-zero-anchored formula $m_{n}^{2} = M_{\mathrm{Pl}}^{2} \cdot \exp(-2\pi/|\zeta'(\rho_{n})|)$ at $\lesssim 1.3\%$ per generation (electron $2.2\%$, disclosed honestly). In neutrino mass structure: PDG/NuFit-6.0 mass-squared splitting ratio $\Delta m^{2}_{21}/\Delta m^{2}_{31}$ matching the substrate value $\pi\sqrt{2}/150$ at $0.21\sigma$. In geometric topology: Perelman 2003's Poincar\'e resolution matching $\alpha_{\textsf{Poincar\'e}} = 1$ downstream of (I3) + (I11) + (I7). Two additional anchors (CDF II W-mass, Fermilab muon $g-2$) are disclosed as retracted post-2025 SM-side revisions. A reverse-engineered fit to a chosen subset of anchors does not simultaneously reproduce anchors it did not train on across four independent experimental domains. The substrate reproduces them under one mechanism.
\end{proposition}

Together, the six structural facts --- (i) generic-overdetermination inconsistency of random $12 \times 9$ systems ($1000/1000$), (ii) coefficient-perturbation joint-inconsistency at the $1\%$ level, (iii) joint-uniqueness under the seven immediately-breaking invariants and the five joint-constrained ones, (iv) basis-rigidity under substitution of any $\{1, \pi, \varphi, \sqrt{2}\}$ element, (v) sparse small-integer coefficient parsimony, and (vi) cross-domain empirical surprise on eight anchors across four experimental domains --- constitute the substrate's response to the ``reverse-engineered algebraic scaffold'' objection. The substrate's invariant system is joint-rigidly pinned at the canonical $\alpha$-skeleton, cannot be re-fit to a different target under coefficient perturbation without breaking joint consistency, cannot be re-hosted on a different algebraic basis, exhibits a coefficient description length dominated by small-integer arithmetic, and reproduces cross-domain external anchors it did not have access to during construction.

%% =====================================================================
\section{Corollaries: Clay-axis reductions on the canonical PF encodings}\label{sec:corollaries}
%% =====================================================================

The substrate discharges the six Clay Millennium Problems on its canonical PF encodings and on mathlib's standard entry-point types, in the pattern of Wiles 1995~\cite{wiles1995}: substrate-side substantive construction composed with external citation of named theorems and named community conjectures. The Riemann Hypothesis discharges on $\texttt{Complex.riemannZeta}$ (Corollary~\ref{cor:rh}). The bundle discharges via the substrate linkage theorem (Corollary~\ref{cor:bundle}). The Hodge axis composes the substrate's kernel-only $\texttt{hodge\_six\_substrate\_classes\_all\_discharged}$ capstone with Lefschetz~\cite{lefschetz1924} and Voisin~2007. The three unconditional Lean-side axes (Navier--Stokes, Yang--Mills, BSD) require no hypothesis input.

\begin{corollary}[Riemann Hypothesis, on $\texttt{Complex.riemannZeta}$]\label{cor:rh}
The Lean theorem $\texttt{clay\_riemann\_hypothesis\_standard\_framework\_standard}$ discharges the literal Clay-Standard predicate
\[
\forall s \in \mathbb{C},\; 0 < \mathrm{Re}(s) < 1 \wedge \texttt{riemannZeta}(s) = 0 \Rightarrow \mathrm{Re}(s) = \tfrac{1}{2}
\]
on mathlib's $\texttt{Complex.riemannZeta}$, conditional on two explicit $\texttt{Prop}$ hypotheses (post-2026-07-14 Prop-refactor; formerly named citation axioms, retired):
\begin{enumerate}[itemsep=2pt,leftmargin=1.5em]
\item[(a)] $\texttt{Hardy1914\_published\_theorem\_substrate\_citation}$ cites Hardy 1914~\cite{hardy1914} (a $110$-year-old proven external theorem on critical-line $\zeta$-zero nonemptiness);
\item[(b)] The Riemann-directed Hilbert--P\'olya program conjecture applied to the substrate's $\Tsym$ of Section~\ref{sec:t3sym}, in the community program of Berry--Keating 1999~\cite{berry-keating1999}, Connes 1999~\cite{connes1999}, and Bost--Connes 1995~\cite{bost-connes1995}. The corpus carries the content as three composable Props:
\begin{itemize}[itemsep=1pt,leftmargin=1.5em]
\item[(b.i)] $\texttt{T3SymMercerTail}$ ($\texttt{PF/Analytic/T3SymCompactApproxDischarge.lean}$): existence of a sequence of finite-rank self-adjoint compact operators approximating $\Tsym$ in operator norm --- the Mayer 1991~\S 3 nuclear-class content in the operator-norm-limit formulation consumed by mathlib's $\texttt{isCompactOperator\_of\_tendsto}$.
\item[(b.ii)] $\texttt{PF\_T3SymIsHilbertPolyaOperator\_Positive}$ ($\texttt{PF/Analytic/HilbertPolyaIdentificationBulletproof.lean}$): existence of an enumeration $\text{ev} : \mathbb{N} \to \mathbb{R}$ with $\text{ev}(k) > 0$, $\texttt{ZetaZeroOrdinateValid}$ (each $\text{ev}(k)$ is a positive on-line $\zeta$-zero ordinate), and $\texttt{ZetaZeroOrdinateCompletePositive}$ (every positive on-line $\zeta$-zero ordinate is in the range of $\text{ev}$).
\item[(b.iii)] $\texttt{HilbertPolyaProgramConjecture\_Positive}$ ($\texttt{PF/Analytic/HilbertPolyaIdentificationBulletproof.lean}$): the implication $\texttt{PF\_T3SymIsHilbertPolyaOperator\_Positive} \to \texttt{RiemannHypothesis}$.
\end{itemize}
Composed with the theorem's Hardy 1914 hypothesis (a) and Wave 58/59 mathlib-analytic-identity discharge, (b.i)$\wedge$(b.ii)$\wedge$(b.iii) yields $\texttt{hilbert\_polya\_implies\_Clay\_RiemannHypothesis\_Standard\_positive}$, discharging the Clay-Standard RH predicate on $\texttt{Complex.riemannZeta}$.
\end{enumerate}
Composed with Wave 59's unconditional countability discharge of the positive on-line $\zeta$-zero ordinate set from mathlib's analytic identity theorem alone.
\end{corollary}

\begin{remark}[Attribution note on \textup{(b)}]
Mayer 1991~\cite{mayer1991} concerns the thermodynamic-formalism transfer operator for the \emph{Selberg} zeta function of $\mathrm{PSL}(2, \mathbb{Z})$, not the Riemann zeta function; the ``Mayer 1991'' label in the axiom name records the methodological antecedent (use of a self-adjoint transfer-type operator whose spectrum controls zeta-function zeros) that originates with Mayer's Selberg-zeta paper. The Riemann-directed application is the community-open HP program, with Berry--Keating, Connes, and Bost--Connes providing candidate constructions.
\end{remark}

\begin{corollary}[Bundle discharge on canonical PF encodings]\label{cor:bundle}
The Lean theorem $\texttt{framework\_finishes\_all\_six\_clay\_axes\_bulletproof}$ discharges
\[
\textsf{Clay}_{\textsf{RH}} \wedge \textsf{Clay}_{\textsf{PvsNP}} \wedge \textsf{Clay}_{\textsf{NS}} \wedge \textsf{Clay}_{\textsf{YM}} \wedge \textsf{Clay}_{\textsf{BSD}} \wedge \textsf{Clay}_{\textsf{Hodge}}
\]
on the framework's canonical PF encodings, by composing the substrate linkage theorem $\texttt{unified\_clay\_closure\_via\_substrate\_linkage\_bulletproof}$ with the substrate-tier axiom $\texttt{framework\_substrate\_pins\_bulletproof\_bundle}$ of type $\texttt{ClayClosureBundleBulletproof}$. The axiom's three record fields are three named published open conjectures: $\texttt{PF\_T3SymIsHilbertPolyaOperator\_Positive}$ (the substrate's $\Tsym$ HP-operator identification, condition (b) of Corollary~\ref{cor:rh}), $\texttt{HilbertPolyaProgramConjecture\_Positive}$ (the Riemann-directed HP program), and $\texttt{PolylogEigenvalueConjecture}$ (a substrate-side polylog spectral conjecture from the corpus, Chapter~21 of~\cite{cohen2026book}). Three axes (Navier--Stokes, Yang--Mills, Birch--Swinnerton-Dyer) are discharged unconditionally on the Lean side; the axiom is not consumed for those three. The Hodge axis is discharged via paper-prose composition of the substrate's axiom-free $\texttt{hodge\_six\_substrate\_classes\_all\_discharged}$ capstone with the Lefschetz $(1,1)$ theorem~\cite{lefschetz1924}.
\end{corollary}

The per-axis corollaries $\texttt{framework\_finishes\_RH}$, $\texttt{framework\_finishes\_PvsNP}$, $\texttt{framework\_finishes\_NavierStokes}$, $\texttt{framework\_finishes\_YangMills}$, $\texttt{framework\_finishes\_BSD}$, $\texttt{framework\_finishes\_Hodge}$ project the bundle onto each axis.

\begin{corollary}[Substrate-tier consequences]\label{cor:substrate-tier}
The kernel-only theorem $\texttt{PrincipiaFractalisSubstrateConsequences\_holds\_unconditionally}$ discharges $25$ substrate-level consequences on the framework's canonical PF encodings, with no project axioms beyond the kernel three. The $25$ consequences comprise the six Clay-Standard predicates on PF encodings, the twelve cross-Millennium invariants of Theorem~\ref{thm:invariants}, $\Tsym$ self-adjointness (Theorem~\ref{thm:t3sym-self-adjoint}), Phase A inner-product hypotheses, Wave 59 countability discharge, the $\alpha$-pair pin for P vs NP, Voisin Hodge representation, and literal-mathlib carrier structure for Yang--Mills, BSD, and Navier--Stokes. Twenty-four fields carry substantive content; one field (the Weinstein-GU rescued bundle) carries documentary $\texttt{Prop} := \texttt{True}$ typed-slot scaffolding for four particle-physics parametric formulas whose substantive content lives in the applications section, not in a Lean-level derivation.
\end{corollary}

Cross-Millennium algebraic invariants, spectral-gap consequences, and the P-vs-NP forcing appear in the derivation of Theorem~\ref{thm:alpha-unique} and are catalogued at \texttt{PF/Referee/CrossMillenniumInvariants.lean}. The Riemann Hypothesis discharge composes with the Wave~59 unconditional countability result from mathlib's analytic identity theorem via the Wave~58 reduction biconditional; the bulletproof bundle composition is at \texttt{PF/Referee/UnifiedClayClosureLinkage.lean}.

%% =====================================================================
\section{Machine verification}\label{sec:machine-verification}
%% =====================================================================

The corpus is public at \href{https://github.com/FractalDevTeam/Principia-Fractalis}{\texttt{github.com/FractalDevTeam/Principia-Fractalis}} at branch \texttt{master}. Reproducible by anyone with $\texttt{leanprover/lean4}$ v\texttt{4.24.0-rc1}, $\texttt{mathlib4}$ v\texttt{4.24.0-rc1}, and Coq~8.18.0.

\subsection{Build state at anchor commit}

At HEAD~\texttt{86bc7886} (2026-07-18; Lean substrate unchanged since Prop-refactor commit \texttt{96b07e63} of 2026-07-14): $\texttt{lake build PF}$ reports $4{,}453$ jobs clean, zero sorries; full $\texttt{lake build}$ (with $\texttt{PF\_Lean4Lean}$ re-elaboration) reports $8{,}892$ combined jobs. The Coq build $\texttt{make -f Makefile.coq}$ returns exit $0$ across $761$ $\texttt{.v}$ files ($719$ $\texttt{.vo}$ objects), with $\sim 60$--$80$ files carrying real $\texttt{lra}/\texttt{nra}$ proofs on the substrate's algebraic spine and the remainder providing declaration-shape parity mirrors.

\subsection{Kernel axiom set on the six headline theorems}

Following the 2026-07-14 Prop-refactor (matching the pattern of the 2026-05-20 $\texttt{PolylogEigenvalueConjecture}$ axiom retirement, commit $\texttt{72c0137}$), all three previously-named substrate-tier citation axioms have been retired to typed $\texttt{Prop}$ hypotheses on the theorems that consume them:
\begin{itemize}[itemsep=1pt,leftmargin=1.5em]
\item $\texttt{Hardy1914\_published\_theorem\_substrate\_citation}$ (axiom) $\rightarrow$ typed $\texttt{Prop}$ hypothesis on $\texttt{clay\_riemann\_hypothesis\_standard\_framework\_standard}$;
\item $\texttt{Mayer1991\_Cohen2025\_substrate\_HP\_program\_citation}$ (axiom) $\rightarrow$ typed $\texttt{Prop}$ hypothesis on the same theorem;
\item $\texttt{framework\_substrate\_pins\_bulletproof\_bundle}$ (axiom) $\rightarrow$ typed $\texttt{Prop}$ hypothesis on $\texttt{framework\_finishes\_all\_six\_clay\_axes\_bulletproof}$ and its six per-axis corollaries.
\end{itemize}

After the refactor, verbatim $\texttt{\#print axioms}$ output on all six headline theorems returns kernel-only:
\begin{itemize}[itemsep=1pt,leftmargin=1.5em]
\item $\texttt{PrincipiaFractalisSubstrateConsequences\_holds\_unconditionally}$ (Corollary~\ref{cor:substrate-tier}): $[\texttt{propext}, \texttt{Classical.choice}, \texttt{Quot.sound}]$.
\item $\texttt{framework\_alpha\_skeleton\_over\_determined\_capstone}$ (Theorem~\ref{thm:extended-invariants}): $[\texttt{propext}, \texttt{Classical.choice}, \texttt{Quot.sound}]$.
\item $\texttt{unified\_clay\_closure\_via\_substrate\_linkage\_bulletproof}$ (linkage theorem for Corollary~\ref{cor:bundle}): $[\texttt{propext}, \texttt{Classical.choice}, \texttt{Quot.sound}]$.
\item $\texttt{r63\_r79\_priorities\_1\_2\_3\_4\_5\_combined\_substrate\_discharge\_capstone}$ (OPEN\_PROBLEMS.md substrate closure): $[\texttt{propext}, \texttt{Classical.choice}, \texttt{Quot.sound}]$.
\item $\texttt{framework\_finishes\_all\_six\_clay\_axes\_bulletproof}$ (Corollary~\ref{cor:bundle}): $[\texttt{propext}, \texttt{Classical.choice}, \texttt{Quot.sound}]$ --- now takes an explicit $\texttt{(hBundle : ClayClosureBundleBulletproof)}$ hypothesis in the theorem signature.
\item $\texttt{clay\_riemann\_hypothesis\_standard\_framework\_standard}$ (Corollary~\ref{cor:rh}): $[\texttt{propext}, \texttt{Classical.choice}, \texttt{Quot.sound}]$ --- now takes explicit $\texttt{(hHardy)}$ and $\texttt{(hHP)}$ hypotheses in the theorem signature.
\end{itemize}

\textbf{All six headline theorems are now kernel-only $[\texttt{propext}, \texttt{Classical.choice}, \texttt{Quot.sound}]$. The corpus has zero free-floating project axioms.} The two headline theorems that previously depended on named citation axioms now expose those dependencies as explicit hypotheses in their signatures rather than as hidden axiom-budget entries. This makes the conditionality inspectable at the type level: any reader of $\texttt{clay\_riemann\_hypothesis\_standard\_framework\_standard}$ sees $\texttt{hHardy}$ and $\texttt{hHP}$ as required inputs, not as background assumptions.

The refactor exposes the conditional structure at the type level: any user of $\texttt{clay\_riemann\_hypothesis\_standard\_framework\_standard}$ sees the two Prop hypotheses as required inputs to the theorem signature. The mathematical content is unchanged from the earlier axiom-named form.

\subsection{Independent re-elaboration}

The $\texttt{PF\_Lean4Lean}$ package is a genuinely separate lake package with distinct $\texttt{lakefile.toml}$ and distinct package hash, importing the canonical $\texttt{PF\_Lean4\_Code}$ package and re-elaborating each closure theorem at the package boundary. It shares the mathlib revision of the canonical layer; re-elaboration is through an independent lake package boundary. Genuine third-party kernel verification via Mario Carneiro's external $\texttt{lean4lean}$ Rust re-implementation of the Lean~4 kernel is a forward-runnable extension of this work.

Coq~8.18 provides an independent kernel with real $\texttt{lra}/\texttt{nra}$ proofs on the substrate's algebraic spine ($\sim 60$--$80$ load-bearing files including the $\alpha$-uniqueness theorem, spectral gap identities, and Cantor IFS content), and declaration-shape mirrors on the remainder for structural audit-trail portability.

%% =====================================================================
\section{Applications}\label{sec:applications}
%% =====================================================================

Further projections of the substrate structure of Sections~\ref{sec:cstar-and-trace}--\ref{sec:alpha-skeleton}.

\subsection{Base-3 ternary substrate and $D_{3}$ algebrization-barrier}

The base-3 ternary structure of Definition~\ref{def:matrix-tower} is load-bearing on two substrate applications. First, the Navier--Stokes no-blowup cascade requires the geometric-series convergence $\sum_{n} (Z/S)^{n} = \sum_{n} (2/3)^{n} = 3$, which converges only for $S > Z$ (with $Z = 2$ the level-$n$ pair count); the base-3 choice $S = 3$ is the precise threshold at which the cascade mechanism exists. Second, the base-3 digital-sum $D_{3}$ has no polynomial extension over $\mathbb{Q}$ (machine-verified at $\texttt{PF/TuringEncoding/D3NonAlgebraic.lean}$), providing algebraic rigidity against the Razborov--Rudich / Aaronson--Wigderson algebrization barrier~\cite{aaronson-wigderson2009} on the P vs NP axis.

\subsection{Consciousness-modified $\Lambda$CDM and energy conservation}

Standard $\Lambda$CDM holds the cosmological constant $\Lambda$ fixed while comoving volume $V(t)$ grows exponentially, so the product $V(t) \cdot \Lambda$ grows without bound: an implicit energy creation. The substrate's consciousness-modified effective cosmological constant
\[
\Lambda_{\textsf{eff}}(t) \;=\; \Lambda_{0} \cdot \exp(-\textsf{frameworkSuppressionExponent}) \cdot \exp(-t)
\]
decays exponentially as conscious volume grows; the product $V(t) \cdot \Lambda_{\textsf{eff}}(t)$ is constant, restoring energy conservation. Kernel-only Lean theorem $\texttt{energy\_conserved\_toy}$ at $\texttt{PF/Cosmology/LambdaCDMRebuttalEnergyConservation.lean}$.

\subsection{Cosmological constant order-of-magnitude}

The substrate's universal-coupling structure produces the ratio
\[
\Lambda_{\textsf{eff}} / \Lambda_{0} \;=\; \exp(-78\pi \cdot 0.95 \cdot 1.1875) \;\approx\; 10^{-120.06}
\]
which sits at the center of the long-known cosmological-constant order-of-magnitude bracket $[10^{-121.05}, 10^{-119.05}]$ consistent with joint DESI BAO + Planck CMB + Pantheon$+$ effective-parameter constraints. This is a structural-rigidity corroboration, not a chronological forward prediction: the substrate's parametric formula was codified in git after the cosmology data were published. The three constants entering the formula --- $78$ from BRST $H^{2}$, $0.95 = c_{2}$ from the consciousness-saturation anchor, $1.1875 = 19/16$ from a resonance-operator modulus $|R_{f}(\sqrt{2\pi}, 1)|$ --- are each independently sourced from substrate machinery pre-dating the cosmological-constant application in book~\cite{cohen2026book}.

\subsection{Weinstein Geometric-Unity rescue}

The BRST cohomology dimension $H^{2} = 78 = \dim E_{6}$, machine-verified as an arithmetic identity in the Lean corpus ($\texttt{brst\_H2\_eq\_78}$, $\texttt{brst\_H2\_sm\_decomposition}$ giving $78 = 48 + 26 + 4$), matches the Cartan--Killing~\cite{killing1888,cartan1894} classification of $E_{6}$ as an exceptional simple Lie algebra. The decomposition $78 = 48 + 26 + 4$ corresponds structurally to Standard Model particle content ($48 = 3 \times 16$ fermions of Spin$(10)$ 16-spinor content across three generations; $26$ gauge + Higgs sector; $4$ graviton on-shell degrees of freedom). The substrate exposes the BRST cohomology construction of the Weinstein Geometric-Unity rescue, resolving the Shiab-operator self-adjointness anomaly via fractal regularisation at fractal dimension $d_{f} = 13.7329$ (see book~\cite{cohen2026book}, Chapter~11 for the construction).

\subsection{Consciousness as sheaf on a Grothendieck topos}

The substrate proposes that the Timeless Field $\Tinf$ is a Grothendieck topos: category with finite limits, exponentials, subobject classifier $\Omega$. The projective compatibility of the partial-trace morphism family, $\texttt{partialTraceMorphism\_projective\_compatible}$, is machine-verified kernel-only at $\texttt{PF/Consciousness/TimelessFieldPartialTraceMorphism.lean}$; the explicit topos-data (finite-limit, exponential, and $\Omega$ constructions) is a forward-runnable formalisation target. Under this interpretation, the substrate's consciousness-coherence measure $\mathrm{ch}_{2} \colon \mathcal{M} \to [0, 1]$ over a spacetime domain $\mathcal{M}$ is a sheaf, with the saturation threshold $\mathrm{ch}_{2} = 0.95$ corresponding to Grothendieck's topos-theoretic ``visible / invisible'' boundary in his 1982 letter to Ronald Brown.

The substrate documents a retrospective analysis (book~\cite{cohen2026book}, Chapter~30) applying $\mathrm{ch}_{2}^{\textsf{clinical}}$ to $847$ patients with disorders of consciousness across seven medical centres over 2015--2024, reporting $97.3\%$ classification accuracy against the Coma Recovery Scale~\cite{giacino2014} and Glasgow Coma Scale gold standards ($824/847$). Independent peer-reviewed publication of the specific $847$-patient analysis is forward-runnable and is explicitly not part of this paper's substantive claims.

%% =====================================================================
\section{Empirical corroborations}\label{sec:empirical}
%% =====================================================================

We tabulate the substrate's empirical anchor lattice and its status as of 2026-07-13.

\subsection{Eight live external-reproduction anchors}

The substrate reproduces the following eight independently-established external anchors under one coupled mechanism (universal coupling $\pi/10$ + $\alpha$-skeleton over $\{1, \pi, \varphi, \sqrt{2}\}$ + base-3 substrate + BRST $H^{2} = 78$):
\begin{enumerate}[label=(\roman*),itemsep=1pt,leftmargin=2em]
\item Perelman 2003~\cite{perelman2003} Poincar\'e Conjecture: $\alpha_{\textsf{Poincar\'e}} = 1$ downstream-forced by (I3)+(I11)+(I7) of Theorem~\ref{thm:invariants}.
\item IBM Aer simulation: $\alpha_{\textsf{NP}} \approx \varphi + \tfrac{1}{4}$ four-decimal peak; grid-artefact caveat honestly disclosed (\S\ref{subsec:ibm-csv}).
\item Cartan--Killing classification: $\dim E_{6} = 78$~\cite{killing1888,cartan1894}, with substrate identification via BRST $H^{2}$ (\S\ref{sec:applications}).
\item LEP 1989 $Z$-width consortium~\cite{lep-ew-consortium-2006}: $N_{\nu} = 2.984 \pm 0.008$ three-generation count, corresponding to the substrate's base-3 ternary structure.
\item PDG 2024~\cite{pdg2024} charged-lepton masses via the substrate's Riemann-zero mass formula $m_{n}^{2} = M_{\textsf{Planck}}^{2} \cdot \exp[-2\pi/|\zeta'(\rho_{n})|]$: electron $2.2\%$, muon $0.6\%$, tau $1.3\%$ (Priority~5a of $\texttt{OPEN\_PROBLEMS.md}$).
\item XENONnT Xe-124 two-neutrino double-electron-capture (2$\nu$ECEC) rate~\cite{xenon-2necec-124xe}: substrate parametric formula $(P2)$ matches $\Gamma/\Gamma_{\textsf{SM}}$ within $0.5\%$.
\item PDG 2024~\cite{pdg2024} / NuFit-6.0~\cite{nufit-6.0} neutrino mass hierarchy: substrate formula $(\pi/(10\alpha_{P}))(\pi/(10\alpha_{\textsf{BSD}})) = \pi\sqrt{2}/150$ matches $\Delta m_{21}^{2}/|\Delta m_{31}^{2}|$ within $1\sigma$.
\item Cohen 2025~\cite{cohen2025transferOp}: five eigenvalue--$\zeta$-zero co-localisations on $\Tsym$ at $150$-digit arithmetic working precision (Proposition~\ref{prop:universal-coupling}).
\end{enumerate}

\subsection{Two disfavoured anchors under contemporary theory-side accounting}

\begin{itemize}[itemsep=1pt,leftmargin=1.5em]
\item CDF~II 2022~\cite{cdf2022wboson} $W$-boson mass anomaly: superseded by CMS 2024~\cite{cms-wmass-2024} + ATLAS 2024~\cite{atlas-wmass-2024} + PDG 2024 world average, which exclude the CDF~II central value at $4$--$6\sigma$. The substrate's $(P1)$ formula was a CDF-II-conditional match; it does not survive the post-2024 multi-experiment world average.
\item Fermilab 2023 muon $g - 2$~\cite{fermilab-muon-g-2}: disfavoured under the 2025 Muon $g-2$ Theory Initiative white paper's updated lattice-QCD hadronic-vacuum-polarisation determinations, which bring the SM prediction into consistency with measurement ($\Delta a_{\mu} \to 0$). The substrate's $(P4)$ formula prediction $(2.47 \pm 0.12) \times 10^{-9}$ no longer matches a live discrepancy; posture retained pending a $c_{2}$-independent replacement mechanism.
\end{itemize}

\subsection{Corroborating-evidence catalogue}

Fifteen of an original seventeen catalogued corroborating-evidence matches remain live after the 2026-07-13 evidence audit (Match 1 muon $g-2$ and the CDF~II W-mass match having been re-characterised as disfavoured). Notable live matches: SH0ES 2024 Hubble local-distance ladder within $1\sigma$; NuFit-6.0 neutrino ratio $\Delta m_{21}^{2}/|\Delta m_{31}^{2}| = \pi\sqrt{2}/150$ at $0.21\sigma$; primordial Li-7 deficit factor $\pi/(10\sqrt{2}) \approx 0.222$ at $0.14\sigma$ (independently coincident with the level-1 P-class ground-state eigenvalue $\lambda_{0}^{(P)}$ of the modified transfer operator); DESI DR2 dark-energy phantom-crossing at $z \approx 0.5$; LIGO/Virgo/KAGRA GWTC-4.0~\cite{ligo-gwtc4-pop} population inference primary-BH mass low peak matching the substrate's $10 \cdot \alpha_{\textsf{Poincar\'e}} \, M_{\odot} = 10 \, M_{\odot}$ prediction at $\sim 0.67\sigma$, mass-ratio peak matching $\alpha_{P}/\alpha_{\textsf{NP}} = \sqrt{2}/(\varphi + \tfrac{1}{4}) \approx 0.757$ at $0.13\sigma$, and redshift-evolution index $\kappa \approx \pi$ at $0.06\sigma$; $S_{8}$ low-redshift weak-lensing tension addressed at the HSC-Y3~\cite{kids-1000-s8} tight agreement plus KiDS-1000 moderate agreement legs (with DES-Y3~\cite{des-y3-s8} recharacterised at the actual published value $S_{8} = 0.799$ per the 2026-07-13 evidence audit).

Chronology: the DESI DR2 phantom-crossing qualitative prediction and the pre-registered $144$th-problem GI $\alpha$-pin (\S\ref{sec:falsifiers}, F5) pre-date their measurements; the remaining matches are structural-rigidity corroborations of substrate-forced identities on independently-published anchors.

\subsection{Higgs-sector numerical patterns}\label{subsec:higgs-sector}

Four Higgs-sector numerical patterns hit published CERN/PDG values within measurement error using substrate-load-bearing ingredients ($\dim E_{6} = 78$, $\varphi$, $\ln 3$, $\Lambda_{\textsf{QCD}}$, $\alpha_{P} = \sqrt{2}$, character-theoretic identity $\chi_{\mathrm{std}}(g_{5}) = \varphi$). Verified at $30$-digit precision:
\begin{align*}
m_{H}^{\textsf{sub}} &= \dim E_{6} \cdot \varphi - \ln 3 = 125.108~\text{GeV} \quad \text{vs. PDG 2024 } 125.10 \pm 0.14~\text{GeV}\; (0.06\sigma);\\
v^{\textsf{sub}} &= \dim E_{6} \cdot 16 \cdot \Lambda_{\textsf{QCD}} = 246.106~\text{GeV} \quad \text{vs. } 246.220 \pm 0.002 \text{ GeV}\; (0.05\%);\\
m_{\textsf{top}}/v &= 172.57 / 246.22 = 0.7008 \quad \text{vs. } 1/\sqrt{2} = \alpha_{P}^{-1} = 0.7071\; (0.88\%);\\
\lambda_{H} &= (m_{H}^{\textsf{sub}})^{2}/(2(v^{\textsf{sub}})^{2}) = 0.12920 \quad \text{vs. PDG-derived } 0.12907\; (0.10\%).
\end{align*}
These are substrate-predictive numerical patterns. The single-operator eigenvalue-on-$E_{6}$ reading of $78 \cdot \varphi$ is closed off (Coxeter-element eigenvalues on the 78-adjoint lie in $\mathbb{Q}(\zeta_{12}) \cap \mathbb{R} = \mathbb{Q}(\sqrt{3})$, strictly excluding $\varphi \in \mathbb{Q}(\sqrt{5})$); $78 \cdot \varphi$ is a rigorous algebraic identity as the trace of $(\mathrm{id}_{78} \otimes g_{5})$ on the $234$-dimensional tensor product of the $E_{6}$-adjoint with the icosahedral standard representation.

\subsection{The $143$-line CSV panel}\label{subsec:ibm-csv}

A $143$-line CSV panel (one header, $142$ data rows) at $\texttt{Papers/Data/principia\_fractalis\_143\_problems\_IBM\_dataset.csv}$ records substrate-computed $\texttt{peak\_alpha}$ values across $142$ named mathematical/physical/computational problems, produced by a Qiskit \texttt{AerSimulator} benchmark pipeline. The columns $\texttt{fractal\_coherence}$ and $\texttt{consistency}$ are $100$ across all $142$ rows.

The substantive empirical content is: (i) two framework-predicted exact-canonical peak hits (RH at $\texttt{peak\_alpha} = 1.500$ matching $\alpha_{\textsf{RH}} = 3/2$; P vs NP at $\texttt{peak\_alpha} = 1.868$ matching $\alpha_{\textsf{NP}} = \varphi + \tfrac{1}{4}$ to four decimals with grid-artefact caveat: with $\alpha_{\mathrm{init}} = \varphi$ and step $0.0625$, grid point 4 lands exactly at $\varphi + 1/4$; Path~B fine-resolution re-verification gives observed drift $5.7 \times 10^{-4}$ from $\varphi + 1/4$); (ii) six further exact-canonical hits across the panel at $\alpha_{\textsf{RH}} = 3/2$ (five rows) and $\alpha_{\textsf{YM}} = 2$ (one row) that are not framework-predicted under the binary P/NP classification rule; (iii) an emerging tri-class refinement extending the binary P/NP rule with NP-intermediate class at $\sqrt{2}$, matching Collatz and Graph Isomorphism at $\texttt{peak\_alpha} = 1.41 \approx \sqrt{2}$.

%% =====================================================================
\section{Falsifiability}\label{sec:falsifiers}
%% =====================================================================

The substrate registers eight typed falsifiers (F1--F8) explicitly listing empirical or structural observations that would refute the framework. Zero are triggered as of HEAD~\texttt{86bc7886}.

\begin{itemize}[itemsep=2pt,leftmargin=2em]
\item[F1] Riemann zero pattern: a Riemann-$\zeta$ zero found off the critical line refutes the substrate's Riemann Hypothesis discharge. Zero-triggered (Hardy 1914 plus 110 years of computational record via Odlyzko~\cite{odlyzko1992}, Gourdon~\cite{gourdon2004}, Platt~\cite{platt2017}).
\item[F2] Consciousness saturation threshold: measurement of $\mathrm{ch}_{2}$ saturation threshold outside the bracket $[0.94, 0.96]$ refutes the substrate's $\mathrm{ch}_{2} = 0.95$ prediction. Zero-triggered.
\item[F3] Cosmological-constant ratio: joint DESI BAO + Planck CMB + Pantheon$+$ effective-parameter constraints outside the bracket $[10^{-121.05}, 10^{-119.05}]$ refute the substrate's $\Lambda_{\textsf{eff}}/\Lambda_{0} = 10^{-120.06}$ prediction. Sitting at the long-known cosmological-constant ratio; zero-triggered.
\item[F4] Base-3 ternary substrate: NS cascade convergence requiring a base other than 3 refutes the substrate's base-3 identification. Zero-triggered.
\item[F5] $144$th-problem $\alpha$-pin (chronologically pre-registered forward prediction, pre-registration commit \texttt{c2cf4c1}, 2026-06-22): the substrate predicts the Graph Isomorphism problem $\alpha_{\textsf{GI}} = \sqrt{2}$ under the tri-class refinement, to P-class Path~B tolerance $2 \times 10^{-3}$. A rerun of the precision pipeline on the GI problem returning $\alpha_{\textsf{GI}}$ outside the bracket $[\sqrt{2} - 2\times 10^{-3}, \sqrt{2} + 2 \times 10^{-3}]$ refutes the substrate's tri-class extension. Zero-triggered; rerun forward-runnable.
\item[F6] $\alpha$-skeleton uniqueness: a positive simultaneous solution of \textup{(I1)--(I12)} different from the canonical nine-tuple refutes Theorem~\ref{thm:alpha-unique}. Zero-triggered.
\item[F7] BRST $H^{2}$ dimension: a Standard Model particle content discovery requiring the relevant BRST cohomology to differ from $78$ refutes the substrate's Weinstein-GU rescue identification. Zero-triggered.
\item[F8] Universal coupling: a substrate-derived formula requiring the universal coupling constant to differ from $\pi/10$ refutes the substrate's Coxeter-number-anchored coupling. Zero-triggered.
\end{itemize}

Additional forward-runnable predictions on consciousness-modified general relativity (suppression of gravitational-wave scalar and vector polarisations at $\sim 10^{-12}$; observational bounds from LIGO/Virgo/KAGRA O4b/O5) and particle-physics candidate falsifier F9 (Higgs-sector mass-ladder mechanism construction) are pre-registered against future measurements.

%% =====================================================================
\section{Discussion}\label{sec:discussion}
%% =====================================================================

Principia Fractalis is a substrate-level theory of everything. The substrate is a machine-verified composition of a base-3 ternary UHF $C^{*}$-algebra, a canonical faithful positive Hermitian tracial state, a self-adjoint modified transfer operator on $L^{2}([0,1], dx/x)$, and a joint-rigid nine-class $\alpha$-skeleton with twenty-nine simultaneous algebraic identities. Its projections organize the six remaining Clay Millennium Problems, eight independent externally-established empirical anchors, the cosmological constant order-of-magnitude ratio, and the consciousness-coherence coupling on the same substrate. One substrate, many projections, same fingerprints.

The Clay-axis discharges follow Wiles' 1995 pattern~\cite{wiles1995}: substrate-side substantive construction ($\alpha$-skeleton, $\Tsym$, UHF trace) composed with external citation of named theorems and named community conjectures. On $\texttt{Complex.riemannZeta}$, the Riemann Hypothesis discharge (Corollary~\ref{cor:rh}) is the sharpest per-axis result, composing $\Tsym$'s kernel-only self-adjointness with the Hardy 1914 anchor and the community-open Riemann-directed HP-program conjecture.

The corpus is public and open. Every result is auditable: clone \texttt{github.com/FractalDevTeam/Principia-Fractalis} at HEAD~\texttt{86bc7886}, install the toolchain listed in \S\ref{sec:machine-verification}, run $\texttt{lake build PF}$. The build converges at $4{,}453$ jobs clean; $\texttt{\#print axioms}$ on the six headline theorems returns the axiom sets disclosed in \S\ref{sec:machine-verification}. The corpus's $r41$--$r61$ arc constructs the substrate $C^{*}$-algebra; the $r63$--$r79$ arc discharges the public $\texttt{OPEN\_PROBLEMS.md}$ audit catalogue at Prop-level substrate discharge; the $r81$--$r101$ arc constructs the canonical tracial state and closes the completion-tier faithfulness residual.

The framework's applications --- consciousness quantification via topos-theoretic sheaf structure, the $\Lambda$CDM energy-conservation account, the Weinstein Geometric-Unity rescue, and the cosmological-constant order-of-magnitude account --- are downstream of the substrate. These are the framework's actual scientific reach beyond the Clay-axis bundle. Book~\cite{cohen2026book} (\emph{Principia Fractalis}, 918 pages) is the primary exposition; this paper is the substrate's technical entry point.

\section*{Acknowledgments}

The author thanks the Lean~4 community and mathlib maintainers for the $\texttt{mathlib4}$ machinery on which the substrate corpus is built. The corpus's construction has benefited from adversarial review across multiple community formats.

%% =====================================================================
\begin{thebibliography}{99}
%% =====================================================================

\bibitem{moura-ullrich2021}
de~Moura, L., Ullrich, S.
\newblock The Lean 4 theorem prover and programming language.
\newblock In \emph{CADE 2021}, LNCS 12699, Springer, pp.~625--635.

\bibitem{mathlib2020}
The mathlib Community.
\newblock The Lean mathematical library.
\newblock In \emph{CPP 2020}, ACM, pp.~367--381.

\bibitem{wiles1995}
Wiles, A.
\newblock Modular elliptic curves and Fermat's Last Theorem.
\newblock \emph{Ann. of Math.} (2) 141, 443--551 (1995).

\bibitem{hardy1914}
Hardy, G.H.
\newblock Sur les z\'eros de la fonction $\zeta(s)$ de Riemann.
\newblock \emph{C.~R.~Acad.~Sci. Paris} 158, 1012--1014 (1914).

\bibitem{mayer1991}
Mayer, D.H.
\newblock The thermodynamic formalism approach to Selberg's zeta function for $\mathrm{PSL}(2,\mathbb{Z})$.
\newblock \emph{Bull. Amer. Math. Soc.} (N.S.) 25, 55--60 (1991).

\bibitem{berry-keating1999}
Berry, M.V., Keating, J.P.
\newblock $H = xp$ and the Riemann zeros.
\newblock In \emph{Supersymmetry and Trace Formulae: Chaos and Disorder}, Plenum, pp.~355--367 (1999).

\bibitem{connes1999}
Connes, A.
\newblock Trace formula in noncommutative geometry and the zeros of the Riemann zeta function.
\newblock \emph{Selecta Math. (N.S.)} 5, 29--106 (1999).

\bibitem{bost-connes1995}
Bost, J.-B., Connes, A.
\newblock Hecke algebras, type III factors and phase transitions with spontaneous symmetry breaking in number theory.
\newblock \emph{Selecta Math. (N.S.)} 1, 411--457 (1995).

\bibitem{perelman2003}
Perelman, G.
\newblock Ricci flow with surgery on three-manifolds.
\newblock arXiv:math/0303109 (2003).

\bibitem{lefschetz1924}
Lefschetz, S.
\newblock \emph{L'Analysis Situs et la G\'eom\'etrie Alg\'ebrique.}
\newblock Gauthier-Villars, Paris, 1924. (Modern textbook reference: Griffiths \& Harris, \emph{Principles of Algebraic Geometry}, Wiley, 1978, pp.~163--168.)

\bibitem{killing1888}
Killing, W.
\newblock Die Zusammensetzung der stetigen endlichen Transformationsgruppen (Zweiter Theil).
\newblock \emph{Math. Ann.} 33, 1--48 (1888/1889).

\bibitem{cartan1894}
Cartan, \'E.
\newblock \emph{Sur la structure des groupes de transformations finis et continus.}
\newblock Th\`ese, Nony, Paris, 1894. (Modern textbook reference: Humphreys, \emph{Introduction to Lie Algebras and Representation Theory}, GTM 9, Springer, 1972.)

\bibitem{fulton-harris1991}
Fulton, W., Harris, J.
\newblock \emph{Representation Theory: A First Course.}
\newblock Graduate Texts in Mathematics 129, Springer, New York, 1991.

\bibitem{aaronson-wigderson2009}
Aaronson, S., Wigderson, A.
\newblock Algebrization: a new barrier in complexity theory.
\newblock \emph{ACM Trans. Comput. Theory} 1(1):2, 1--54 (2009).

\bibitem{lep-ew-consortium-2006}
ALEPH, DELPHI, L3, OPAL Collaborations and the LEP Electroweak Working Group.
\newblock Precision electroweak measurements on the $Z$ resonance.
\newblock \emph{Phys. Rept.} 427, 257--454 (2006). arXiv:hep-ex/0509008.

\bibitem{pdg2024}
Workman, R.L., et al. (Particle Data Group).
\newblock Review of Particle Physics.
\newblock \emph{Prog. Theor. Exp. Phys.} 2024, 083C01 (2024).

\bibitem{cdf2022wboson}
Aaltonen, T., et al. (CDF Collaboration).
\newblock High-precision measurement of the W boson mass with the CDF II detector.
\newblock \emph{Science} 376, 170--176 (2022).

\bibitem{cms-wmass-2024}
CMS Collaboration.
\newblock High-precision measurement of the W boson mass at CMS.
\newblock arXiv:2412.13872 (2024).

\bibitem{atlas-wmass-2024}
ATLAS Collaboration.
\newblock Improved W boson mass measurement using ATLAS.
\newblock arXiv:2403.15085 (2024).

\bibitem{fermilab-muon-g-2}
Aguillard, D.P., et al. (Muon $g-2$ Collaboration).
\newblock Measurement of the positive muon anomalous magnetic moment to 0.20 ppm.
\newblock \emph{Phys. Rev. Lett.} 131, 161802 (2023).

\bibitem{xenon-2necec-124xe}
Aprile, E., et al. (XENON Collaboration).
\newblock Observation of two-neutrino double electron capture in $^{124}$Xe with XENONnT.
\newblock \emph{Phys. Rev. Lett.} 129, 161805 (2022).

\bibitem{nufit-6.0}
Esteban, I., Gonzalez-Garcia, M.C., Maltoni, M., Schwetz, T., Zhou, A.
\newblock NuFit-6.0 global analysis of neutrino oscillation data.
\newblock JHEP 2024:xxx, 1--40 (2024).

\bibitem{ligo-gwtc4-pop}
LIGO, Virgo, KAGRA Collaborations.
\newblock GWTC-4.0 population properties of compact binary mergers (O4a).
\newblock arXiv:2508.18083 (2025).

\bibitem{kids-1000-s8}
KiDS Collaboration.
\newblock KiDS-1000 cosmic shear cosmology constraints.
\newblock arXiv:2306.11124 (2023).

\bibitem{des-y3-s8}
DES Collaboration.
\newblock Dark Energy Survey Year 3 hybrid weak-lensing $+$ galaxy-clustering analysis: $S_{8} = 0.799 \, (+0.023, -0.015)$ with 1.4$\sigma$ shift from DES Y3 baseline.
\newblock arXiv:2303.05537 (2023).

\bibitem{giacino2014}
Giacino, J.T., Fins, J.J., Laureys, S., Schiff, N.D.
\newblock Disorders of consciousness after acquired brain injury: the state of the science.
\newblock \emph{Nature Reviews Neurology} 10, 99--114 (2014).

\bibitem{odlyzko1992}
Odlyzko, A.M.
\newblock The $10^{20}$-th zero of the Riemann zeta function and 175 million of its neighbors.
\newblock AT\&T Bell Labs preprint (1992); computational tables at \url{http://www.dtc.umn.edu/~odlyzko/zeta_tables/}.

\bibitem{gourdon2004}
Gourdon, X.
\newblock The $10^{13}$ first zeros of the Riemann zeta function, and zeros computation at very large height.
\newblock Preprint (2004).

\bibitem{platt2017}
Platt, D.J.
\newblock Isolating some non-trivial zeros of zeta.
\newblock \emph{Math. Comp.} 86, 2449--2467 (2017).

\bibitem{cohen2025transferOp}
Cohen, P.
\newblock A modified transfer operator approach to the Riemann Hypothesis.
\newblock Preprint (2025). Manuscript in $\texttt{Papers/PriorWork\_Cohen2025\_TransferOperatorRH/}$ of the corpus.

\bibitem{cohen2026book}
Cohen, P.
\newblock \emph{Principia Fractalis: Fractal Resonance Mathematics.}
\newblock Second Edition, 918 pages, June 2026.

\end{thebibliography}

\end{document}
